

\vspace{-2mm}
\section{Experiments}
\label{sec:exp}
In this section, we evaluate \mot on two challenge TKGQA datasets. The detailed experimental settings, including datasets, baselines, and implementations, can be found in Appendix \ref{appen:dataset_details}.

% \begin{table}[]
% \caption{Performance comparison of different models
% (in percentage) on MultiTQ}
% \label{tab:main_results_multiTQ}
% \resizebox{1\columnwidth}{!}{

% \begin{tabular}{ccccccc}

% \hline
%                                     &                                 & \multicolumn{5}{c}{{\color[HTML]{0563C1} {\ul \textbf{Hit@1}}}}                                                            \\ \cline{3-7} 
%                                     &                                 &                                    & \multicolumn{2}{c}{\textbf{Question Type}} & \multicolumn{2}{c}{\textbf{Answer Type}} \\ \cline{4-7} 
% \multirow{-3}{*}{\textbf{Model}}    & \multirow{-3}{*}{\textbf{LLM}}  & \multirow{-2}{*}{\textbf{Overall}} & \textbf{Multiple}     & \textbf{Single}    & \textbf{Entiy}      & \textbf{Time}      \\ \hline
% BERT                                &                                 & 8.3                                & 6.1                   & 9.2                & 10.1                & 4.0                \\
% DistillBERT                         &                                 & 8.3                                & 7.4                   & 8.7                & 10.2                & 3.7                \\
% ALBERT                              & \multirow{-3}{*}{-}             & 10.8                               & 8.6                   & 11.6               & 13.9                & 3.2                \\ \hline
% EmbedKGQA                           &                                 & 20.6                               & 13.4                  & 23.5               & 29.0                & 0.1                \\
% CronKGQA                            &                                 & 27.9                               & 13.4                  & 33.7               & 32.8                & 15.6               \\
% MultiQA                             & \multirow{-3}{*}{-}             & 29.3                               & 15.9                  & 34.7               & 34.9                & 15.7               \\ \hline
% ChatGPT                             &                                 & 10.2                               & 7.7                   & 14.7               & 13.7                & 2.0                \\
% KG-RAG                              &                                 & 18.5                               & 16.0                  & 20.0               & 23.0                & 7.0                \\
% CoT KB                              &                                 & 24.0                               & 12.0                  & 44.0               & 22.0                & 32.0               \\
% ReAct KB                            &                                 & 21.0                               & 13.6                  & 63.5               & 31.3                & 30.0               \\
% ARI                                 &                                 & 38.0                               & 68.0                  & 21.0               & 39.4                & 34.4               \\
% TempAgent                           & \multirow{-6}{*}{GPT-3.5-Turbo} & 53.9                               & 16.8                  & 68.4               & 47.8                & 66.1               \\ \hline
%                                     & GPT-4o-mini                     & 64.2                               & 40.3                  & 73.0               & 61.5                & 70.8               \\
%                                     & Qwen3-32B                       & 68.2                               & 42.5                  & 77.6               & 62.1                & 81.7               \\
%                                     & Deepseek-V3                     & {\ul 73.0}                         & {\ul 45.9}            & {\ul 82.9}         & {\ul 67.7}          & {\ul 84.6}         \\
% \multirow{-4}{*}{\textbf{MemoTime}} & GPT-4-Turbo                     & \textbf{77.9}                      & \textbf{53.8}         & \textbf{86.8}      & \textbf{74.5}       & \textbf{85.3}      \\ \hline
% \end{tabular}
% }



% \end{table}

% % Please add the following required packages to your document preamble:
% % \usepackage{multirow}
% % \usepackage[table,xcdraw]{xcolor}
% % Beamer presentation requires \usepackage{colortbl} instead of \usepackage[table,xcdraw]{xcolor}
% % \usepackage[normalem]{ulem}
% % \useunder{\uline}{\ul}{}
% \begin{table}[]

% \end{table}

% \begin{table}[t]
% \caption{Performance comparison of different models
% (in percentage) on TimeQuestions} 
% \label{tab:main_results_timeQ}
% \resizebox{1\columnwidth}{!}{

% \begin{tabular}{ccccccc}
% \toprule
%                                     &                                & \multicolumn{5}{c}{  \textbf{Hit@1}}                                                           \\ \cline{3-7} 
%                                     &                                &                                    & \multicolumn{2}{c}{\textbf{Question Type}} & \multicolumn{2}{c}{\textbf{Answer Type}} \\ \cline{4-7} 
% \multirow{-3}{*}{\textbf{Model}}    & \multirow{-3}{*}{\textbf{LLM}} & \multirow{-2}{*}{\textbf{Overall}} & \textbf{Explicit}    & \textbf{Implicit}   & \textbf{Temporal}   & \textbf{Ordinal}   \\ \midrule
% PullNet                             &                                & 10.5                               & 2.2                  & 8.1                 & 23.4                & 2.9                \\
% Uniqorn                             &                                & 33.1                               & 31.8                 & 31.6                & 39.2                & 20.2               \\
% GRAFT-Net                           & \multirow{-3}{*}{-}            & 45.2                               & 44.5                 & 42.8                & 51.5                & 32.2               \\ \hline
% CronKGQA                            &                                & 46.2                               & 46.6                 & 44.5                & 51.1                & 36.9               \\
% TempoQR                             &                                & 41.6                               & 46.5                 & 3.6                 & 40.0                & 34.9               \\
% EXAQT                               &                                & 57.2                               & 56.8                 & 51.2                & 64.2                & 42.0               \\
% TwiRGCN                             &                                & 60.5                               & 60.2                 & 58.6                & 64.1                & 51.8               \\
% LGQA                                & \multirow{-5}{*}{-}            & 52.9                               & 53.2                 & 50.6                & 60.5                & 40.2               \\ \hline
% GenTKGQA                            & Fine-tune           & 58.4                               & 59.6                 & {\ul 61.1}          & 56.3                & {\ul 57.8}         \\
% TimeR$^4$                                  & Fine-tune  & 64.8                               & {\ul 66.0}           & 52.9                & \textbf{77.6}       & 45.5               \\
% ChatGPT                             & GPT-3.5-Turbo                        & 45.9                               & 43.3                 & 51.1                & 46.5                & 48.1               \\ \hline
%                                     & GPT-4o-mini                        & 60.9                               & 60.0                 & 50.0                & 69.4                & 52.8               \\
%                                     & Qwen3-32B                      & 60.6                               & 58.3                 & 51.1                & 69.8                & 53.3               \\
%                                     & Deepseek-V3                    & {\ul 67.80}                        & 63.1                 & 53.3                & 71.3                & 54.3               \\
% \multirow{-4}{*}{\textbf{MemoTime}} & GPT-3.5-Turbo                        & \textbf{71.4}                      & \textbf{67.0}        & \textbf{67.5}       & {\ul 74.5}          & \textbf{58.7}      \\ \bottomrule
% \end{tabular}
% }
% \end{table}

% % 

% ========== TIMEQUESTIONS ==========
\begin{table}[t]
\centering
\caption{Performance comparison on \textbf{TimeQuestions} (Hits@1, \%). Best in \textbf{bold}, second-best \underline{underlined}.}
\vspace{-3mm}
\label{tab:main_results_timeQ}
\resizebox{\columnwidth}{!}{
\begin{tabular}{l c c c c c c}
\toprule
\multirow{2}{*}{\textbf{Model}} & \multirow{2}{*}{\textbf{LLM}} & \multirow{2}{*}{\textbf{Overall}} & \multicolumn{2}{c}{\textbf{Question Type}} & \multicolumn{2}{c}{\textbf{Answer Type}} \\
\cmidrule(lr){4-5}\cmidrule(lr){6-7}
 &  &  & \textbf{Explicit} & \textbf{Implicit} & \textbf{Temporal} & \textbf{Ordinal} \\
\midrule
PullNet\cite{sun2019pullnet}   & \multirow{3}{*}{--} & 10.5 & 2.2  & 8.1  & 23.4 & 2.9 \\
Uniqorn\cite{pramanik2024uniqorn}   &                      & 33.1 & 31.8 & 31.6 & 39.2 & 20.2 \\
GRAFT-Net\cite{sun2018open} &                      & 45.2 & 44.5 & 42.8 & 51.5 & 32.2 \\
\midrule
CronKGQA\cite{mavromatis2022tempoqr}  & \multirow{5}{*}{--} & 46.2 & 46.6 & 44.5 & 51.1 & 36.9 \\
TempoQR\cite{mavromatis2022tempoqr}   &                      & 41.6 & 46.5 & 3.6  & 40.0 & 34.9 \\
EXAQT\cite{jia2021complex}     &                      & 57.2 & 56.8 & 51.2 & 64.2 & 42.0 \\
TwiRGCN\cite{sharma2022twirgcn}   &                      & 60.5 & 60.2 & 58.6 & 64.1 & 51.8 \\
LGQA\cite{liu2023local}      &                      & 52.9 & 53.2 & 50.6 & 60.5 & 40.2 \\
\midrule
GenTKGQA\cite{gao2024two}  & Fine-tune & 58.4 & 59.6 & \underline{61.1} & 56.3 & \underline{57.8} \\
TimeR$^4$\cite{T4} & Fine-tune & 64.8 & \underline{66.0} & 52.9 & \textbf{77.6} & 45.5 \\
ChatGPT   & GPT-3.5-Turbo & 45.9 & 43.3 & 51.1 & 46.5 & 48.1 \\
\midrule
\multirow{4}{*}{\textbf{MemoTime}}
 & GPT-4o-mini  & 60.9 & 60.0 & 50.0 & 69.4 & 52.8 \\
 & Qwen3-32B    & 60.6 & 58.3 & 51.1 & 69.8 & 53.3 \\
 & DeepSeek-V3  & \underline{67.8} & 63.1 & 53.3 & 71.3 & 54.3 \\
 &\cellcolor{blue!6}  GPT-4-Turbo  & \cellcolor{blue!6} \textbf{71.4} & \cellcolor{blue!6} \textbf{67.0} & \cellcolor{blue!6} \textbf{67.5} & \cellcolor{blue!6} \underline{74.5} & \cellcolor{blue!6} \textbf{58.7} \\
\bottomrule
\end{tabular}
}
% \vspace{-2mm}
\end{table}


\begin{table*}[t]
\centering
\caption{Performance comparison between the IO baseline and \mot across two temporal QA datasets using six backbone LLMs. 
The highest improvement is highlighted in bold, and the second-best is underlined for each data.}
\label{tab:backbone_memotime}
\vspace{-2mm}
\resizebox{1\textwidth}{!}{
\begin{tabular}{l
ccc ccc ccc ccc ccc ccc}
\toprule
\multirow{2}{*}{\textbf{Dataset}} &
\multicolumn{3}{c}{\textbf{Qwen3-4B}} &
\multicolumn{3}{c}{\textbf{Qwen3-8B}} &
\multicolumn{3}{c}{\textbf{Qwen3-32B}} &
\multicolumn{3}{c}{\textbf{Qwen3-80B}} &
\multicolumn{3}{c}{\textbf{DeepSeek-V3}} &
\multicolumn{3}{c}{\textbf{GPT-4-Turbo}} \\[2pt]
\cmidrule(lr){2-4}\cmidrule(lr){5-7}\cmidrule(lr){8-10}\cmidrule(lr){11-13}\cmidrule(lr){14-16}\cmidrule(lr){17-19}
&\textbf{ IO }& \textbf{\mot} & \textbf{$\uparrow$}
& \textbf{IO }& \textbf{\mot }& \textbf{$\uparrow$}
& \textbf{IO} & \textbf{\mot} & \textbf{$\uparrow$}
& \textbf{IO} & \textbf{\mot }& \textbf{$\uparrow$}
& \textbf{IO }& \textbf{\mot} & \textbf{$\uparrow$}
&\textbf{ IO }& \textbf{\mot }& \textbf{$\uparrow$} \\
\midrule
\textbf{MultiTQ }
& 3.5  & 55.3 & 14.8 $\times$ 
&\cellcolor{blue!6} 2.5  &\cellcolor{blue!6} {\ul 57.0} & \cellcolor{blue!6}{\ul 21.8 $\times$ }
& \cellcolor{blue!6}1.3  & \cellcolor{blue!6}\textbf{61.4} &\cellcolor{blue!6} \textbf{46.2 $\times$ }
& 9.5  & 67.5 & 6.1 $\times$ 
& 7.5  & 70.9 & 8.5 $\times$ 
& 12.0 & 76.3 & 5.4 $\times$ \\
\textbf{TimeQuestion}
& \cellcolor{blue!6}18.0 &\cellcolor{blue!6} \textbf{45.3} & \cellcolor{blue!6}\textbf{1.5 $\times$ }
& 28.5 & 49.4 & 0.7 $\times$ 
&\cellcolor{blue!6} 30.0 &\cellcolor{blue!6} {\ul 60.6} & \cellcolor{blue!6}{\ul1.0 $\times$ }
& 41.0 & 62.7 & 0.5 $\times$ 
& 43.5 & 64.6 & 0.5 $\times$ 
& 44.0 & 68.1 & \text{0.5 $\times$ } \\
\bottomrule
\end{tabular}}
% \vspace{-3mm}
\end{table*}



% \vspace{-3mm}
\subsection{Main results}  
% We present the experimental results in comparison
% between our model and the existing state-of-the-art
% baseline models on the MultiTQ and TimeQuestions datasets in Table \ref{tab:main_results_multiTQ} and Table \ref{tab:main_results_timeQ}.

% \myparagraphquestion{Main results on temporal question answering}  
We evaluate \mot against a comprehensive set of baselines on 
% on two of the most challenging temporal benchmarks, 
\text{MultiTQ} and \text{TimeQuestions}, as shown in Table~\ref{tab:main_results_multiTQ} and Table~\ref{tab:main_results_timeQ}. 
% 
% Across both datasets, MemoTime achieves consistent state-of-the-art (SOTA) performance, demonstrating its superiority in temporal reasoning, multi-entity synchronization, and long-term experience reuse.  
On the \text{MultiTQ} dataset, traditional pre-trained language models (BERT, ALBERT) and embedding-based methods (EmbedKGQA, CronKGQA, MultiQA) show limited temporal reasoning ability, with overall accuracies below 30\%.  
Their static embeddings fail to capture multi-hop or cross-entity temporal dependencies.  
LLM-based methods, e.g, KG-RAG, ReAct KB, ARI, TempAgent, leverage prompt reasoning but remain sensitive to implicit time expressions and struggle with multi-constraint temporal alignment.  
In contrast, MemoTime consistently outperforms all competitors across both question and answer types.  
MemoTime with GPT-4-Turbo achieves an overall 77.9\% Hit@1, surpassing TempAgent by 24.0\% and outperforming all GPT-3.5 baselines by a large margin.  
Even smaller backbones (e.g., Qwen3-32B) reach 68.2\%, exceeding the best GPT-3.5-based models.  
These results confirm that MemoTime effectively integrates temporal grounding, hierarchical reasoning, and operator-aware retrieval to ensure temporal faithfulness and reasoning stability.
On the \text{TimeQuestions} dataset, which emphasizes explicit versus implicit temporal expressions and temporal versus ordinal answer types, MemoTime again achieves the best overall results.  
As shown in Table~\ref{tab:main_results_timeQ}, MemoTime with GPT-4-Turbo achieves 71.4\% overall accuracy and 74.5\% on temporal-type questions, outperforming the fine-tuned GenTKGQA and TimeR$^4$ models, despite being fully training-free.  
This demonstrates MemoTime’s ability to handle implicit temporal relations and operator combinations through memory-guided decomposition and dynamic temporal retrieval.  
Overall, across both datasets, MemoTime achieves new state-of-the-art results under all backbone LLMs, validating its design as a \text{memory-augmented temporal reasoning framework} capable of learning from experience, synchronizing multiple entities, and maintaining temporal consistency across diverse question types.


\vspace{-2mm}



\subsection{Ablation Study}\label{exp:ablation}
% We perform various ablation studies to understand the importance of different factors in \mot. 
% The ablation studies are conducted with GPT-3.5-Turbo on two subsets of the test sets of CWQ and WebQSP, each of which contains 500 randomly sampled questions.
% These ablation studies are performed with GPT-3.5-Turbo on two subsets of the CWQ and WebQSP test sets, each containing 500 randomly sampled questions.
% \newline

\myparagraphquestion{How does the performance of \mot vary across different LLM backbones}
To evaluate the generality of our approach, we tested \mot on six LLM backbones, including Qwen3-4B, Qwen3-8B, Qwen3-32B, Qwen3-80B, DeepSeek-V3, and GPT-4-Turbo, across two temporal QA benchmarks (MultiTQ and TimeQuestions). 
As shown in Table~\ref{tab:backbone_memotime}, \mot consistently improves performance over the IO baseline across all settings. 
On the more complex \text{MultiTQ}, which requires multi-hop and cross-entity temporal reasoning, \mot boosts accuracy from as low as 1.3\% to 61.4\% on Qwen3-32B, achieving a \text{46.2 times} relative gain. 
Even the smallest backbones (e.g., Qwen3-4B) experience a \text{14.8 times} improvement, demonstrating that our recursive, memory-augmented reasoning compensates for weaker temporal understanding in smaller models.
For \text{TimeQuestions}, which involves simpler, single-hop temporal relations, \mot still improves performance modestly by up to 1.5 times, indicating consistent stability across reasoning difficulty levels.  
Overall, \mot enables smaller models to perform competitively with large-scale LLMs such as GPT-4-Turbo, while stronger models continue to benefit from enhanced temporal faithfulness and reduced reasoning variance. 
These results confirm that \mot generalizes well across architectures and scales, serving as a plug-in temporal reasoning module rather than relying on model-specific fine-tuning.


\begin{table}[t]
\centering
% \vspace{-3mm}
\caption{Ablation study on \textbf{MultiTQ} (Hits@1, \%).}
\vspace{-2mm}
\label{tab:ablation_memotime}
\resizebox{\columnwidth}{!}{
\begin{tabular}{lccccc}
\toprule
\multirow{2}{*}{\textbf{Model Variant}} & \multirow{2}{*}{\textbf{Overall}} & 
\multicolumn{2}{c}{\textbf{Question Type}} & 
\multicolumn{2}{c}{\textbf{Answer Type}} \\ 
\cmidrule(lr){3-4} \cmidrule(lr){5-6}
& & \textbf{Multiple} & \textbf{Single} & \textbf{Entity} & \textbf{Time} \\
\midrule
\textbf{MemoTime (Full w/GPT-4o-mini)}                & \textbf{64.2} & \textbf{40.3} & \textbf{73.0} & \textbf{61.5} & \textbf{70.8} \\
\textit{w/o Graph-based Retrieval}      & 52.9          & 23.5          & 65.0          & 48.2          & 64.2          \\
\textit{w/o Embedding Retrieval}        & 60.1          & 38.6          & 68.0          & 57.4          & 66.8          \\
\textit{w/o Temporal Evidence Retrieval}& 11.2          & 8.0           & 12.4          & 11.5          & 12.0          \\
\textit{w/o Question Tree}              & 58.3          & 19.3          & 72.6          & 54.6          & 70.1          \\
\textit{w/o Experience Memory}          & 59.8          & 26.1          & 72.2          & 55.8          & 69.6          \\
\bottomrule
\end{tabular}}
% \vspace{-2mm}
\end{table}


\myparagraphquestion{How does temporal evidence retrieval affect performance}
To assess the contribution of the temporal evidence retrieval module, we disable different retrieval components.
As shown in Table~\ref{tab:ablation_memotime}, removing graph-based retrieval reduces the overall accuracy from 64.2\% to 52.9\%, with a substantial drop on multiple-entity questions (from 40.3\% to 23.5\%), indicating that explicit structural traversal is essential for maintaining temporally aligned entity paths.
When embedding-based retrieval is removed, the performance decreases moderately from 64.2\% to 60.1\%, suggesting that semantic retrieval complements the graph search by capturing lexically diverse or contextually implicit temporal cues.
In contrast, when all external retrieval modules are disabled and the model relies solely on its internal knowledge, the accuracy falls drastically to 11.2\%, reflecting a near-complete loss of temporal grounding.
These findings confirm that temporal evidence retrieval is not only indispensable but also synergistic: graph-based retrieval ensures factual precision through topological consistency, while embedding-based retrieval enhances coverage through semantic generalization.






\myparagraphquestion{How does hierarchical decomposition affect performance} 
To examine the role of hierarchical decomposition, we disable the Tree-of-Time module, forcing the model to reason over the full question without recursive planning.
As shown in Table \ref{tab:ablation_memotime},
the overall accuracy increases from 64.2\% to 58.3\%, while performance on multiple-entity questions drops sharply from 40.3\% to 19.3\%.
This substantial decline demonstrates that recursive decomposition is critical for handling complex temporal dependencies and multi-hop constraints.
Even with retrieval retained, the absence of decomposition limits the model’s ability to enforce monotonic timestamp ordering and to synchronize temporal operators across entities, leading to degraded reasoning coherence.

\myparagraphquestion{How does experience memory affect performance} 
To analyse the role of continual learning, we remove the experience memory module, preventing \mot from accessing previously verified reasoning traces.
The result in Table \ref{tab:ablation_memotime} shows it causes a moderate performance drop from 64.2\% to 59.8\%, with larger reductions observed on multi-hop questions, from 40.3\% to 26.1\%.
The results indicate that experience memory enhances reasoning stability and efficiency.
Although the model remains competitive without it, the memory mechanism contributes to more consistent reasoning trajectories and progressive self-improvement across inference cycles.

To further evaluate the effectiveness and efficiency of \mot, we conduct additional experiments, including multi-granular temporal QA analysis, efficiency analysis, 
% and interpretability-focused case studies illustrating temporal faithfulness 
and case studies of temporal interpretability and
faithful reasoning
in Appendix~\ref{appendxi:add_exp}.




% To further evaluate the effectiveness and efficiency of \mot, we conduct additional experiments, including multi-granular temporal QA analysis, efficiency analysis, case studies of temporal interpretability and
% faithful reasoning in Appendix~\ref{appendxi:add_exp}, 
% % error analysis (Appendix~\ref{appendix:exp:error_analysis}), 
% % and interpretability-focused case studies illustrating temporal faithfulness 
% specialized temporal retrieval toolkits in Appendix~\ref{sec:toolkit}.
% % and 
% % (Appendix~\ref{sec:casestudy-temporal}).







% and prompt setting ablation (Appendix \ref{appendix:Ablation}), reasoning faithfulness analysis (Appendix \ref{appendix:exp:reasoning_faithful}), error analysis (Appendix \ref{appendix:exp:error_analysis}), LLM calls cost and running time analysis (Appendix \ref{appendix:exp:llm_calls}), and graph reduction and path pruning case study (Appendix \ref{appendix:exp_case_study}). 



% \subsection{Effectiveness Evaluation}

% \myparagraphquestion{How well does MemoTime handle multi-granular temporal reasoning}  
% We further evaluate \mot on temporal reasoning tasks that require distinguishing and aligning facts across different temporal granularities.  
% Traditional pre-trained and embedding-based models exhibit substantial performance degradation as the temporal granularity becomes finer. 
% For instance, while MultiQA achieves 44.5\% at the \text{Equal-Day} level, its accuracy drops sharply to 28.3\% for \text{Equal-Multi-Year}.  
% In contrast, \mot consistently delivers strong performance across all granularities and operators, achieving up to \text{94.1\%} Hits@1 on \text{Before/After-Year} and maintaining stable accuracy even under compound constraints (\text{70.4\%} on \text{Equal-Multi-Month}).  
% These results confirm that MemoTime not only maintains temporal faithfulness at fine-grained levels but also generalizes to mixed-granularity reasoning, demonstrating that MemoTime’s operator-aware decomposition and temporal path construction effectively preserve monotonic timestamp ordering and enable unified reasoning across heterogeneous resolutions.


% \subsection{Analysis}


% \begin{table}
% \caption{Computational cost on \text{MultiTQ} by question types.}
% \label{tab:multitq_cost}
% \centering
% \resizebox{1\linewidth}{!}{
% \begin{tabular}{lccccccc}
% \toprule
% \textbf{} & \textbf{Overall} & \textbf{After-First} & \textbf{Before-Last} & \textbf{Equal-Multi} & \textbf{First/Last} & \textbf{Before/After} & \textbf{Equal} \\
% \midrule
% \textbf{Ave. Depth }     & 1.37 & 2.12 & 2.00 & 1.57 & 1.00 & 1.54 & 1.04 \\
% \textbf{Ave. Branch}     & 2.64 & 3.91 & 3.94 & 3.05 & 2.00 & 2.78 & 2.01 \\
% \textbf{Ave. LLM calls}  & 6.14 & 7.50 & 8.80 & 8.00 & 6.00 & 6.60 & 5.00 \\
% \bottomrule
% \end{tabular}}
% \end{table}
% \myparagraph{Efficiency analysis on MultiTQ}
% \label{exp:efficiency_multitq}
% We analyze the computational efficiency of \mot on \textsc{MultiTQ} by reporting the average reasoning depth, branch factor, and LLM calls per question (Table~\ref{tab:multitq_cost}). 
% Overall, \mot exhibits a shallow reasoning structure (depth 1.37, branch 2.64) with only 6.1 calls per query, reflecting efficient recursive control. 
% Simple operators ( e.g., \text{equal}, \text{first/last}) require minimal recursion, while compound and boundary-sensitive operators (\text{after–first}, \text{before–last}) demand deeper trees and broader search. 
% Mixed operators remain moderate but involve additional sub-branch exploration for temporal alignment. 
% These results show \mot dynamically adapts its decomposition and pruning strategies according to operator complexity, achieving fine-grained temporal reasoning with controlled computational overhead.

